\scanpage{90}
\mspage{2.}
\margin{$i:X_0\to\unsure{\mathfrak X|\mathscr X|\mathcal X}$}

\uline{Lemme 2} Soit $\beta\in H^2(X_0,\underline\mu)$ la classe self-intersection de $X_0$ dans $\mathfrak X$, i.e. \inserted{la classe de \deleted{de} Chern \deleted{\unsure{et}} de} $\deleted{\unsure{\mathcal L}}T_{X_0/\mathfrak X})\simeq i^*(\underline L(X_0))$. Soit $\sigma:H^2(x,\mu)\unsure{-|\to}H^2(U,I)$ l'hom canonique. Alors
\[
 \rho(\alpha)=\sigma(\beta).
\]

Cela détermine l'\uline{extension} $\alpha$, sous réserve de connaître le quotient $I$ de $\mu$. On sait déjà que le \unsure{composé}
\diagram{\begin{tikzcd}[ampersand replacement=\&,column sep=.5em,row sep=.5em]
 {}\&H^2(\widetilde U^t,\widetilde\mu)\simeq\operatorname{Hom}(\pi_2^t(\unsure{X|U}),\mu_0)\&{}\&{}\\
 H^2(U,\mu)\arrow[r]\arrow[ur]\&H^2(U,I)\arrow[r]\&H^2(\widetilde U^t,I_0)=\operatorname{Hom}(\pi_2^t(U),I_0\&{}\\
 \deleted{\unsure{\rho(\beta)}}\quad\beta|U\arrow[u,equal]\&{}\&{}\&{}
\end{tikzcd}}
s'annule sur $\rho(\alpha)$. On a le diagramme (\unsure{comm.|commutatif}\inserted{par lemme 2}),
% slanted
\margin{Du moins si les $D'_i\subset X_0$ ($1\leq i\leq\unsure{n|r}$) sont \uline{disjoints} [à vérifier si cela est indispensable]}

\diagram{\begin{tikzcd}[ampersand replacement=\&,column sep=1em,row sep=1em]
 \pi_2^t(U)=H_2(\widetilde U)\arrow[d,"\overline\beta"]\arrow[dr,"\overline{\rho(\alpha)}"]\arrow[r,dash,"\deleted{\unsure{\longrightarrow}}"]\&\deleted{\pi_2^t(U)=H_2(\widetilde U^t)}\arrow[d,"\alpha","\deleted{\unsure{\rho}}"' ]\&{}\\
 \mu_0\arrow[r]\&\deleted{\unsure{\mu_0}}\unsure{I_0}\arrow[r,dash,"\deleted{\unsure{\longrightarrow}}"]\&\deleted{\unsure{I}}
\end{tikzcd}}
et par lemme 1. Ceci \unsure{bien} \unsure{---} défini l'hom. \deleted{dans} du $\pi_2^t$ défini par une 2-classe ].

Donc \deleted{\unsure{Im}} $\mu_0\to I_0$ s'\uline{annule} \uline{sur} $\operatorname{Im}\pi_2^t(U)$. Je dis que \deleted{\unsure{elle est}} le noyau
% slanted
\margin{\uline{\unsure{Heuristique}}}
est \uline{égal} à cette image ! [En fait, si $I'_0\ \unsure{\mathrm{de}|{}}=\deleted{\unsure{I}}\mu_{\unsure{\infty0|0}}/\operatorname{Im}\pi_2^t(U)$, on aura que $\operatorname{Im}\beta\in\operatorname{Ker}(H^2(U,I'_0)\unsure{-|\to}H^2(\widetilde U^t,I'_0))\simeq H^2(\pi_1^t,I'_0)$, donc $I\beta$ \uline{définit une extension}
% archive sheet number: 89
